\chapter{\texorpdfstring{El Rígido Sincronizado (China): Coerenza di
Comando $\neq$
\(\Phi\)}{El Rígido Sincronizado (China): Coerenza di Comando $\neq$ \textbackslash Phi}}\label{il-rigido-sincronizzato-cina-coerenza-di-comando-phi}

\begin{quote}
\textbf{Argumento central:} la Repubblica Popolare Cinese del 2026 è il
caso paradigmatico di una \textbf{falsa \(\Phi\)}: il sistema esibisce
coerenza visibile massima a livello di comando (\(K \to \infty\) per
imposizione gerarchica, fasi \(\theta_j\) allineate per costrizione
esterna anziché per sincronización spontanea), ma ciò \emph{non è}
l'ordine di Kuramoto. Il parámetro de orden vero è eroso da tre processi
convergenti --- Neijuan estructural (\(\Omega \to \infty\)), colapso
demografico (riduzione catastrofica del sostrato di \(N\) futuro),
decoupling tecnologico forzato (auto-acoplamiento incompleto) ---
mientras la pressione su Taiwan misura il punto di rottura della rigidità.
\end{quote}



\section{Definición Estructural del Rigido
Sincronizzato}\label{definición-estructural-del-rigido-sincronizzato}

\Hypoth{} \textbf{Definizione (Rígido Sincronizado DEQ$^2$).} Un attore
\(e\) è detto \emph{Rígido Sincronizado} se soddisfa simultaneamente:

\begin{enumerate}
\def\labelenumi{\arabic{enumi}.}
\tightlist
\item
  \(K_e \to \infty\) effettivo per imposizione gerarchica (la
  coordinazione delle fasi \(\theta_j\) avviene per \emph{coercizione
  esterna} e non per acoplamiento di Kuramoto autentico);
\item
  \textbf{\(\Phi_e^{\text{osservato}}\) alto,
  \(\Phi_e^{\text{spontaneo}}\) basso} --- esiste un divario estructural
  fra l'ordine \emph{misurato dall'esterno} (apparenza unanimistica) e
  l'ordine \emph{che il sistema produrrebbe in assenza di forzatura};
\item
  \(\Omega_e \uparrow\) in modo monotonico (l'energia dissipata in
  mantenimento dell'allineamento e in conflitto interno latente cresce
  più rapidamente dell'energia utile);
\item
  \(A_e\) \textbf{biforcato}: alto sull'asse materiale-tecnologico
  (capacità di mobilitazione enorme), basso sull'asse
  politico-discorsivo (zero opzionalità di dissenso, por lo tanto zero
  apprendimento dagli errori sistemici);
\item
  \(s_{ii,e}\) alto, \(s_{iii,e}\) in calo \emph{all'esterno} (cresce
  internamente per saturación del campo informativo).
\end{enumerate}

\Proved{} \textbf{Conse siguenza formal.} \(T_s^{(2),\text{sys}}\) è alta
sotto misura interna (numeratore gonfiato da
\(\Phi^{\text{osservato}}\)), ma fragile sotto perturbazione esogena:
una caduta anche piccola di \(K\) (rallentamento della macchina di
forzatura) produce un crollo non lineare di \(\Phi^{\text{osservato}}\)
verso \(\Phi^{\text{spontaneo}}\). Il sistema è \textbf{catastrofico nel
senso di Thom}: la sua transición de fase sarà discontinua, non
graduale come quella dell'Hegemón Decoherente.



\section{Il IV Plenum 2025: Misura della Coerenza
Apparente}\label{il-iv-plenum-2025-misura-della-coerenza-apparente}

\Proved{} \textbf{Dato.} Il IV Plenum del 20$^\circ$ Comitato Centrale del PCC
si è tenuto dal 20 al 23 ottobre 2025. Xi Jinping ha
\emph{personalmente} diretto il drafting team delle raccomandazioni per
il 15$^\circ$ Piano Quinquennale (2026-2030). Il plenum ha riconfermato il
``ruolo del nucleo'' () di Xi senza eccezioni rilevate pubblicamente. Il
piano è centrato su:

\begin{itemize}
\tightlist
\item
  \emph{High-Quality Development} (alta qualità contro alta velocità);
\item
  Auto-sufficienza tecnologica come imperativo economico centrale;
\item
  \emph{Sicurezza economica} come priorità di pari rango con la
  prosperità.
\end{itemize}

\Hypoth{} \textbf{Lettura DEQ$^2$.} Il IV Plenum \emph{misura}
simultaneamente \(K^{\text{imposto}}\) (massimo) e
\(\Phi^{\text{spontaneo}}\) (non osservabile direttamente, ma
\emph{inferibile} dalla congiuntura di tre fatti contemporanei).

{\def\LTcaptype{none} % do not increment counter
{\small\begin{longtable}[]{@{}
  >{\raggedright\arraybackslash}p{(\linewidth - 2\tabcolsep) * \real{0.5000}}
  >{\raggedright\arraybackslash}p{(\linewidth - 2\tabcolsep) * \real{0.5000}}@{}}
\toprule\noalign{}
\begin{minipage}[b]{\linewidth}\raggedright
Fatto del periodo
\end{minipage} & \begin{minipage}[b]{\linewidth}\raggedright
Lettura DEQ$^2$
\end{minipage} \\
\midrule\noalign{}
\endhead
\bottomrule\noalign{}
\endlastfoot
Riconferma del ``core'' Xi senza dissenso pubblico &
\(\Phi^{\text{osservato}} \to 1\) \\
Target GDP 2026 al \(4{,}5\)-\(5\%\) (il più basso da decenni) &
\(\Phi^{\text{spontaneo}}\) già in contrazione: il sistema non riesce
più a \emph{produrre} il proprio ordine economico \\
Lancio della \emph{anti-involution campaign} governativa contro la
sovracapacità produttiva & Riconoscimento \emph{istituzionale} di
\(\Omega \uparrow\) \\
\end{longtable}}
}

\Proved{} \textbf{Conse siguenza.} L'apparente massimo di \(\Phi\) del IV
Plenum è da leggere insieme al riconoscimento ufficiale del Neijuan: il
sistema si è imposto di \emph{parlare di sé} come sistema in difficoltà
estructural. Questa è una firma DEQ$^2$ inequivocabile di Rigido
Sincronizzato --- l'unanimità è ancora completa, ma il contenuto
dell'unanimità è la diagnosi del proprio colapso.



\section{\texorpdfstring{Neijuan (\(\Omega\)) come Sostanza
Strutturale}{Neijuan (\textbackslash Omega) come Sostanza Strutturale}}\label{neijuan-omega-come-sostanza-estructural}

\Hypoth{} \textbf{Tesi.} Il \emph{Neijuan} (, \emph{involución}) ---
concetto antropologico ridefinito da \textbf{Xiang Biao} come \emph{``un
ciclo morto in cui le persone si forzano costantemente, una corsa in cui
ai partecipanti non è permesso fallire né uscire''} --- è la
realizzazione fenomenologica di \(\Omega\) a scala sociale. Non è
ornamento descrittivo: è la \emph{firma empirica diretta} della
variable DEQ$^2$ nel sistema chino.

\Proved{} \textbf{Verifica empirica della crescita di
\(\Omega^{\text{Cina}}\).}

{\def\LTcaptype{none} % do not increment counter
{\small\begin{longtable}[]{@{}
  >{\raggedright\arraybackslash}p{(\linewidth - 6\tabcolsep) * \real{0.2500}}
  >{\raggedright\arraybackslash}p{(\linewidth - 6\tabcolsep) * \real{0.2500}}
  >{\raggedright\arraybackslash}p{(\linewidth - 6\tabcolsep) * \real{0.2500}}
  >{\raggedright\arraybackslash}p{(\linewidth - 6\tabcolsep) * \real{0.2500}}@{}}
\toprule\noalign{}
\begin{minipage}[b]{\linewidth}\raggedright
Indicatore
\end{minipage} & \begin{minipage}[b]{\linewidth}\raggedright
2024
\end{minipage} & \begin{minipage}[b]{\linewidth}\raggedright
2025-2026
\end{minipage} & \begin{minipage}[b]{\linewidth}\raggedright
Lettura
\end{minipage} \\
\midrule\noalign{}
\endhead
\bottomrule\noalign{}
\endlastfoot
Trimestri consecutivi di deflazione & \(\approx 7\) & \textbf{\(10\)}
(al Q1 2026) & dissipazione produttiva estructural \\
PPI in territorio negativo & \(> 2\) anni & \textbf{\(> 3\) anni} &
sovracapacità non riassorbita \\
Property: anno di contrazione & \(4^\circ\) & \textbf{\(5^\circ\)} consecutivo &
crollo del moltiplicatore di settore-chiave (30\% PIL) \\
Property investment (var. annua) & \(-10\%\) ca. &
\textbf{\(-11{,}1\%\)} (feb 2026) & accelerazione del declino \\
GDP target ufficiale & \(5\%\) & \textbf{\(4{,}5\)-\(5\%\)} &
riconoscimento istituzionale del rallentamento \\
Anti-involution campaign & latente & \textbf{dichiarata da Xi} &
\(\Omega\) riconosciuto come problema politico \\
\end{longtable}}
}

\Hypoth{} \textbf{Interpretazione formal.} Il pattern \emph{deflazione
persistente \(+\) sovracapacità \(+\) campagna anti-involutiva} equivale
a \(\dot\Omega > 0\) e \(\ddot\Omega > 0\) simultaneamente. La China ha
prodotto più capacità di quanta il proprio sistema interno possa
assorbire al prezzo che mantiene viva la rete produttiva --- produzione
marginale a costo crescente, donde il rapporto \emph{energia dissipata /
energia utile} cresce monotonicamente.

\Conj{} \textbf{Sotto-tesi}. L'esistenza stessa di una \emph{campagna
ufficiale contro l'involución} è il momento DEQ$^2$ più rivelatore: in un
sistema con \(\Phi\) autentica, \(\Omega\) verrebbe corretto da
meccanismi distribuiti (mercato, dibattito pubblico, federalismo). In un
Rígido Sincronizado, \(\Omega\) può essere corretto solo per
\emph{altra imposizione gerarchica} --- il che ne aumenta a sua volta
\(\Omega\). La campagna anti-involutiva è essa stessa un atto
involutivo.



\section{Decoupling Tecnologico: Auto-acoplamiento
Incompleto}\label{decoupling-tecnologico-auto-acoplamiento-incompleto}

\Proved{} \textbf{Dato di sistema.} Aprile 2026: DeepSeek pubblica V4
con \emph{day-zero adaptation} sui chip Huawei Ascend \(950\)PR e
\(950\)DT. Migrazione da CUDA (Nvidia) a CANN (Huawei) completa: la China
dispone, per la primero volta dal 2022, di uno \emph{stack AI end-to-end
domestico}, dal silicio al modello. Alibaba, ByteDance, Tencent ordinano
centinaia di migliaia di chip Ascend.

\Hypoth{} \textbf{Lettura DEQ$^2$.} La risposta chino al regime di export
controls USA (avviato 2022, intensificato 2025 con il ban H20 --- costo
\(5{,}5\)B charge per Nvidia, vedi Cap.~7.4) configura un
\textbf{auto-acoplamiento forzato}: il sistema si chiude su se stesso
per non dipendere da un esterno percepito come ostile. Da un punto di
vista DEQ$^2$, questo è \(K^{\text{interno}} \uparrow\) (gli accoppiamenti
fra sotto-attori cinesi si rinforzano), ma a costo di:

\begin{enumerate}
\def\labelenumi{\arabic{enumi}.}
\tightlist
\item
  \textbf{Riduzione del campo \(\Pi\) globale} in cui il sistema chino
  è inserito (perdita di varietà \(\Sigma\) --- eredità Topografia);
\item
  \textbf{Aumento di \(A^{\text{tecnologico}}\)} (resiliencia alla
  volatilità da sanzioni);
\item
  \textbf{Riduzione di \(A^{\text{politico-discorsivo}}\)} (perdita di
  interfaccia per recepire segnali esterni dissonanti);
\item
  \textbf{Effetto retroattivo su \(\Phi^{\text{globale}}\)}: se sia USA
  che China si chiudono in stack tecnologici disgiunti, il sistema
  globale frammenta in due cluster con bassa \(\Phi\) inter-cluster.
\end{enumerate}

\Proved{} \textbf{Conse siguenza per \(T_s^{(2),\text{sys}}\) chino.}
L'auto-acoplamiento aumenta numericamente \(\Phi^{\text{osservato}}\)
interno ma \emph{non corregge} \(\Omega\): una macchina che gira più
velocemente su se stessa non riassorbe la sovracapacità, la riproduce
con sostegno tecnologico domestico. La firma è coerente con la diagnosi
\$\S\$8.2.

\Conj{} \textbf{Predizione.} Se la migrazione tecnologica avrà successo
nel triennio 2026-2028, la \emph{visibilità esterna} di
\(\Phi^{\text{Cina}}\) aumenterà (narrativa di vittoria); ma \(\Omega\)
continuerà a crescere porque il modello produttivo che alimenta il
decoupling è lo stesso che produce sovracapacità. Punto di tensione:
l'export è stato l'unico contributo positivo al GDP 2025 (\(+1{,}7\) pp
su crescita aggregata, \emph{più della metà}) --- la chiusura del
mercato USA via dazi (Cap.~7.5) attacca esattamente questa valvola di
sfogo.



\section{\texorpdfstring{Collasso Demografico: l'Erosione del Sostrato
di \(\Phi\)
Futuro}{Collasso Demografico: l'Erosione del Sostrato di \textbackslash Phi Futuro}}\label{colapso-demografico-lerosione-del-sostrato-di-phi-futuro}

\Proved{} \textbf{Dato 2025-2026.} La China registra un \emph{birth rate}
di \(5{,}63\) per \(1000\) persone nel 2025, \textbf{il minimo storico
dal 1949}. Le nascite del 2025 sono il \(-17\%\) con respecto al 2024,
malgrado anni di programmi incentivanti (incluso il sussidio nazionale
del luglio 2025 da 90 miliardi di yuan, \(\approx 12{,}5\)B\$).
Popolazione 2026: \(1{,}413\)B (vs \(1{,}416\)B 2025, perdita di
\(\approx 3{,}2\)M in un anno). Quota over-60: \(23\%\) (2025), \(22\%\)
(2024). Età mediana \(40{,}6\); proiezione 2050: \(52\) anni (mediana
mondiale).

\Hypoth{} \textbf{Lettura DEQ$^2$ formal.} Tutti i parámetros del modello
sono definiti su una popolazione \(N\) di sotto-attori. Il \emph{Rigido
Sincronizzato} la cui \(N\) collassa subisce due effetti distinti:

{\def\LTcaptype{none} % do not increment counter
{\small\begin{longtable}[]{@{}
  >{\raggedright\arraybackslash}p{(\linewidth - 4\tabcolsep) * \real{0.3333}}
  >{\raggedright\arraybackslash}p{(\linewidth - 4\tabcolsep) * \real{0.3333}}
  >{\raggedright\arraybackslash}p{(\linewidth - 4\tabcolsep) * \real{0.3333}}@{}}
\toprule\noalign{}
\begin{minipage}[b]{\linewidth}\raggedright
Effetto
\end{minipage} & \begin{minipage}[b]{\linewidth}\raggedright
Meccanismo
\end{minipage} & \begin{minipage}[b]{\linewidth}\raggedright
Conse siguenza
\end{minipage} \\
\midrule\noalign{}
\endhead
\bottomrule\noalign{}
\endlastfoot
\(\Sigma\) del sistema si contrae & Meno individui \(\Rightarrow\) meno
varianza generazionale e occupazionale &
\(A^{\text{adattivo}} \downarrow\) \\
\(\langle\Omega\rangle / \Phi\) esplode & Il rapporto
pensionati/forza-lavoro \(\to 1\) entro 2050 & \(T_s^{(2),\text{sys}}\)
denominatore \(\to \infty\) \\
\end{longtable}}
}

\Proved{} \textbf{Conse siguenza.} Il colapso demografico è \textbf{non
recuperabile} entro l'orizzonte 2031 della DEQ$^2$: i nati del 2026
entreranno nella forza lavoro después il 2046. La transición de fase
chino è por lo tanto \emph{parámetroszzata da una variable lenta che ha già
attraversato la sua umbral crítico}. Il sistema si sta muovendo verso il
proprio \(T_s^{(2),\text{sys}} \to 0\) con una traiettoria
deterministica del fattore demografico, indipendiente da qualunque scelta
politica corrente.

\Conj{} \textbf{Caveat.} Il fallimento RAND (2025) della politica
pronatalista non è un'anomalia: è coerente con la previsione DEQ$^2$
secondo cui un Rígido Sincronizado non può modificare una variable
lenta come la fertilità con strumenti gerarchici, porque la fertilità
dipende da \(Q_0\) del campo discorsivo (eredità Topografia \S{}5.2 ---
calidad ética, dimensioni \(x, y, z\) del progetto di vita). La leva del
Rígido Sincronizado agisce su \(K^{\text{esterno}}\), non su \(Q_0\).



\section{Taiwan: Banco di Prova della
Rigidità}\label{taiwan-banco-di-prova-della-rigidituxe0}

\Proved{} \textbf{Dato 2024-2026.} Lai Ching-te eletto a gennaio 2024
con il \(40\%\) (DPP perde la maggioranza in Yuan legislativo).
Insediato il 20 maggio 2024. Il 13 marzo 2025 caratterizza la RPC come
\emph{``foreign hostile force''} (forza ostile straniera) --- primo
presidente taiwanese a usare questa formula. Le incursioni PLA in ADIZ
Taiwan: media \(> 300/mese\) después maggio 2024 (più del doppio del
biennio precedente). Esercitazioni \emph{Justice Mission 2025} (29-30
dicembre 2025) --- seconda esercitazione di blocco nel solo 2025.

\Hypoth{} \textbf{Lettura DEQ$^2$.} Taiwan è il punto in cui
\(K^{\text{esterno}}\) chino raggiunge il limite della sua
applicabilità. Tutte le altre variables (Hong Kong, mercato interno,
Tibet, Xinjiang) sono state portate sotto \(K^{\text{imposto}}\) massimo
--- Taiwan no, porque è esterno al perimetro di sovranità. La risposta
del sistema è di \emph{aumentare la pressione coercitiva} (incursioni,
esercitazioni) per simulare \(K^{\text{esterno}}\) anche fuori dai
propri confini.

\Hypoth{} \textbf{Tesi formal.} L'escalation militare chino su Taiwan
misura il \emph{grado di consapevolezza interna del proprio colapso
\(\Phi^{\text{spontaneo}}\)}. Quanto maggiore è il divario fra
\(\Phi^{\text{osservato}}\) e \(\Phi^{\text{spontaneo}}\), tanto più
necessaria diventa una vittoria nazionalistica visibile per ricondurre
temporaneamente i due valori. Taiwan è la \emph{valvola simbolica}
a través de cui il Rígido Sincronizado può tentare di ricostituire
\(\Phi^{\text{osservato}}\) a basso costo materiale immediato --- al
prezzo di una transición de fase potenzialmente catastrofica del
sistema globale.

\Conj{} \textbf{Predizione DEQ$^2$ (datata).} Probabilità di un evento di
tipo \emph{bloqueo} (blocco navale e aereo non escalato in conflitto
generale) nel triennio 2027-2029: \(0{,}25\)-\(0{,}40\) condicional al
peggioramento simultaneo di tre indicatori cinesi (deflazione
persistente \(> 12\) trimestri, GDP \(< 4\%\), ulteriore calo nascite
oltre soglia \(5/1000\)). La firma DEQ$^2$ del Rígido Sincronizado è che
la pressione \emph{aumenta} quando \(\Omega\) aumenta, porque la
coerenza visibile va difesa anche quando --- e specialmente quando ---
quella spontanea si contrae.



\section{\texorpdfstring{BRICS+ e l'Esportazione di Ordine senza
\(\Phi\)}{BRICS+ e l'Esportazione di Ordine senza \textbackslash Phi}}\label{brics-e-lesportazione-di-ordine-senza-phi}

\Proved{} \textbf{Dato.} BRICS espansa nel 2025 con Egitto, Etiopia,
Iran, EAU, Arabia Saudita, Indonesia. Blocco esteso: \(45\%\)
popolazione mondiale, \(35\%\) GDP globale. China-Rusia:
\(\approx 90\%\) scambi bilaterali in valute locali. Yuan share
intra-BRICS: \(\approx 50\%\). Indonesia (23 novembre 2025) lancia
operazioni FX centrate sullo yuan. BRICS Pay operativo target 2030
(slittato da scadenze precedenti). Pilot da 100 unità di token digitale
a copertura mista (40\% oro + 60\% basket BRICS), ottobre 2025.
Presidenza BRICS 2026: India (CBDC interconnection in agenda, non
yuan-centrico).

\Hypoth{} \textbf{Lettura DEQ$^2$.} L'iniziativa chino sui BRICS+ è un
\emph{tentativo di costruire \(\Phi\) esogena} --- esportare l'ordine
all'esterno per compensare la sua erosione interna. Strutturalmente, è
equivalente a una \emph{strategia di diversione di \(K\)}: applicare
l'acoplamiento gerarchico a un sistema più ampio (multilaterale) per
acquisire legitimidad esterna sostitutiva di \(s_{iii}\) in calo.

\Proved{} \textbf{Limite intrinseco.} Un Rígido Sincronizado \emph{non
può generare \(\Phi\) vera} in un sistema multilaterale per ragioni
strutturali: la \(\Phi\) di Kuramoto richiede sotto-attori con fasi
\(\theta_j\) libere che si sincronizzano spontaneamente sotto
\(K > K_c\). Imporre \(K\) esterno (via BRICS Pay, via yuan settlement)
produce \emph{correlazione}, non \emph{coherencia de fase}. Tre indizi
convergenti:

{\def\LTcaptype{none} % do not increment counter
{\small\begin{longtable}[]{@{}
  >{\raggedright\arraybackslash}p{(\linewidth - 2\tabcolsep) * \real{0.5000}}
  >{\raggedright\arraybackslash}p{(\linewidth - 2\tabcolsep) * \real{0.5000}}@{}}
\toprule\noalign{}
\begin{minipage}[b]{\linewidth}\raggedright
Fatto 2025-2026
\end{minipage} & \begin{minipage}[b]{\linewidth}\raggedright
Lettura DEQ$^2$
\end{minipage} \\
\midrule\noalign{}
\endhead
\bottomrule\noalign{}
\endlastfoot
BRICS Pay slitta dal 2025 al 2030 & \(\Phi^{\text{BRICS}}\)
insufficiente per un settlement system reale \\
Pilot token a copertura mista (oro \(+\) basket) & Il sistema \emph{non
si fida del solo yuan} --- China cede sovranità monetaria per acquisire
credibilità multilaterale \\
Presidenza India 2026 spinge CBDC interconnection (non yuan) & India
\emph{resiste} al \(K^{\text{cinese}}\) --- la \(\Phi\) multilaterale
richiede simmetria \\
\end{longtable}}
}

\Conj{} \textbf{Predizione.} La quota dello yuan nelle riserve globali
(IMF COFER) resterà sotto il \(4\%\) nel 2027 e sotto il \(5\%\) nel
2030 --- pur in presenza di dedollarizzazione accelerata (Cap.~7.5). La
dedollarizzazione \emph{non} converge sullo yuan ma su una pluralità di
valute (euro, oro, CBDC, basket BRICS). Falsificabile su dati IMF
trimestrali.



\section{Síntesis del Capítulo}\label{sintesi-del-capítulo}

\Proved{} \textbf{In una frase:} la China del 2026 esibisce coerenza di
comando massima (IV Plenum unanime, ``core Xi'' riconfermato) ma è
strutturalmente afflitta da quattro processi convergenti che erodono il
\(\Phi\) vero --- Neijuan riconosciuto istituzionalmente come problema,
deflazione decennale, colapso demografico irreversibile entro
l'orizzonte DEQ$^2$, decoupling tecnologico riuscito ma a costo di chiusura
sistemica --- e la pressione su Taiwan e BRICS+ misura l'urgenza con cui
il Rígido Sincronizado cerca di sostituire la \(\Phi\) spontanea
perduta con un'imposizione esterna sostitutiva.

{\def\LTcaptype{none} % do not increment counter
{\small\begin{longtable}[]{@{}
  >{\raggedright\arraybackslash}p{(\linewidth - 6\tabcolsep) * \real{0.2500}}
  >{\raggedright\arraybackslash}p{(\linewidth - 6\tabcolsep) * \real{0.2500}}
  >{\raggedright\arraybackslash}p{(\linewidth - 6\tabcolsep) * \real{0.2500}}
  >{\raggedright\arraybackslash}p{(\linewidth - 6\tabcolsep) * \real{0.2500}}@{}}
\toprule\noalign{}
\begin{minipage}[b]{\linewidth}\raggedright
Variabile
\end{minipage} & \begin{minipage}[b]{\linewidth}\raggedright
Valore stimato 2026
\end{minipage} & \begin{minipage}[b]{\linewidth}\raggedright
Direzione
\end{minipage} & \begin{minipage}[b]{\linewidth}\raggedright
Conse siguenza per \(T_s^{(2),\text{sys}}\)
\end{minipage} \\
\midrule\noalign{}
\endhead
\bottomrule\noalign{}
\endlastfoot
\(s_{ii}\) (capacidad material) & molto alta & stabile & mantenuta \\
\(s_{iii}\) (legitimidad esterna) & \(73{,}5\) Brand Finance, in salita
& \(\uparrow\) & rinforzo asimmetrico (vs USA) \\
\(\Phi^{\text{osservato}}\) & \(\approx 1\) (unanimità performativa) &
stabile & numeratore sovrastimato \\
\(\Phi^{\text{spontaneo}}\) & \(\to 0\) (deflazione, Neijuan,
demografia) & \(\downarrow\downarrow\) & numeratore reale eroso \\
\(A^{\text{tecnologico}}\) & alto (DeepSeek-Huawei stack) & \(\uparrow\)
& resiliencia sanzioni \\
\(A^{\text{politico-discorsivo}}\) & basso & \(\downarrow\) & zero
apprendimento da dissonanza \\
\(\Omega\) (Neijuan) & crescente, riconosciuto da Xi &
\(\uparrow\uparrow\) & gonfia denominatore \\
\(N\) futuro (sostrato demografico) & nascite \(5{,}63/1000\) minimo
storico & crollo & \(\Sigma\) contratta entro 2046 \\
\end{longtable}}
}

\Hypoth{} \textbf{Tesi conclusiva del capítulo.}
\(T_s^{(2),\text{sys},\text{Cina}}\) è strutturalmente in regime
\emph{catastrofico} nel senso di Thom: appare alta sotto misura interna,
è in realtà mantenuta da \(K^{\text{imposto}}\) massimo, ma una
qualunque caduta non lineare di \(K\) (rallentamento della macchina di
forzatura --- leadership transition, evento di fiducia, blocco economico
esterno) produce un crollo discontinuo di \(\Phi^{\text{osservato}}\)
verso \(\Phi^{\text{spontaneo}}\). La China è por lo tanto l'attore che,
nell'orizzonte 2028-2031 della DEQ$^2$, ha la \textbf{probabilità più alta
di una transición de fase non graduale}.

\Hypoth{} \textbf{Asimmetria con l'Hegemón Decoherente (Cap.~7).} USA e
China convergono entrambi a \(\Phi \to 0\) ma per percorsi opposti: gli
USA \emph{non hanno mai prodotto} la \(\Phi\) richiesta (decoherencia
distribuita); la China la produce per imposizione e la perderebbe se
l'imposizione cessasse (rigidità catastrofica). Né USA né China producono
\(\Phi\) come \emph{risorsa esportabile} --- ed è qui che si apre lo
spazio del Cap.~9 (Brasil come Bilanciere).



\section{Quattro Predizioni Datate del Cap.~8 (per Cap.
11)}\label{quattro-predizioni-datate-del-cap.-8-per-cap.-11}

\begin{enumerate}
\def\labelenumi{\arabic{enumi}.}
\tightlist
\item
  \textbf{Deflazione chino si estende a \(> 12\) trimestri consecutivi
  entro Q4 2026.} Falsificabile su dati NBS / PIIE / World Bank.
\item
  \textbf{Probabilità evento \emph{bloqueo} su Taiwan nel triennio
  2027-2029: \(0{,}25\)-\(0{,}40\), condicional a deflazione
  persistente \(+\) GDP \(< 4\%\) \(+\) nascite \(< 5/1000\).}
  Falsificabile su corso degli eventi.
\item
  \textbf{Quota yuan in IMF COFER \(< 4\%\) nel 2027 e \(< 5\%\) nel
  2030}, malgrado dedollarizzazione accelerata. Falsificabile su dati
  IMF trimestrali.
\item
  \textbf{Anti-involution campaign chino non riduce \(\Omega\) entro
  2028}: PPI in territorio negativo \(\geq 5\) anni totali (cumulati)
  entro Q4 2027. Falsificabile su dati NBS.
\end{enumerate}



\section{Notas Epistémicas de Cierre
Capítulo}\label{note-epistemiche-di-chiusura-capítulo}

{\def\LTcaptype{none} % do not increment counter
{\small\begin{longtable}[]{@{}
  >{\raggedright\arraybackslash}p{(\linewidth - 4\tabcolsep) * \real{0.3333}}
  >{\raggedright\arraybackslash}p{(\linewidth - 4\tabcolsep) * \real{0.3333}}
  >{\raggedright\arraybackslash}p{(\linewidth - 4\tabcolsep) * \real{0.3333}}@{}}
\toprule\noalign{}
\begin{minipage}[b]{\linewidth}\raggedright
\#
\end{minipage} & \begin{minipage}[b]{\linewidth}\raggedright
Afirmación
\end{minipage} & \begin{minipage}[b]{\linewidth}\raggedright
Estado
\end{minipage} \\
\midrule\noalign{}
\endhead
\bottomrule\noalign{}
\endlastfoot
1 & Definizione formal Rígido Sincronizado (5 condizioni) & \Hypoth{}
categoria DEQ$^2$ nuova \\
2 & IV Plenum 2025 misura \(\Phi^{\text{osservato}} \to 1\) & \Proved{}
dato pubblicato \\
3 & Anti-involution campaign come ammissione istituzionale di \(\Omega\)
& \Hypoth{} inferenza forte ma esplicita nei documenti ufficiali \\
4 & Mappatura Neijuan (Xiang Biao) \(\to \Omega\) DEQ$^2$ & \Hypoth{}
mappatura concettuale forte \\
5 & Stack DeepSeek-Huawei come auto-acoplamiento forzato & \Proved{}
dato aprile 2026 \\
6 & Demografia come variable lenta oltre umbral crítico & \Proved{}
dato 2025-2026 (RAND, Worldometer) \\
7 & Taiwan come valvola di compensazione \(\Phi^{\text{spontaneo}}\) &
\Conj{} hipótesis estructural forte \\
8 & BRICS+ come tentativo di esportare ordine, limite intrinseco &
\Hypoth{} inferenza se derivata \\
9 & Quota yuan COFER \(< 5\%\) 2030 & \Conj{} predizione testabile \\
10 & Probabilità \emph{bloqueo} Taiwan \(0{,}25\)-\(0{,}40\) in
2027-2029 & \Conj{} conjetura datata, falsificabile \\
11 & Asimmetria USA/China (decoherencia distribuita vs rigidità
catastrofica) & \Hypoth{} risultato estructural comparativo \\
\end{longtable}}
}



\section{Fonti dei dati 2025-2026
utilizzate}\label{fonti-dei-dati-2025-2026-utilizzate}

\begin{itemize}
\tightlist
\item
  IV Plenum CCP 20-23 ottobre 2025; piano 2026-2030 --- Bloomberg, CSET
  Georgetown, Chatham House
\item
  Anti-involution campaign --- ufficialmente lanciata da Xi
  (Wikipedia/Neijuan, China Briefing)
\item
  China-Briefing / Rhodium Group / World Bank --- GDP target
  \(4{,}5\)-\(5\%\) 2026, deflazione 10 trimestri
\item
  Property investment \(-11{,}1\%\) feb 2026 --- China Briefing dicembre
  2025
\item
  Xiang Biao (Oxford / Max Planck) --- definición di Neijuan come
  \emph{dead loop}
\item
  DeepSeek V4 + Huawei Ascend \(950\)PR/\(950\)DT --- SCMP, CNN,
  Bloomberg aprile 2026
\item
  Birth rate \(5{,}63/1000\) minimo dal 1949 --- CNBC gennaio 2026,
  Worldometer
\item
  Worldometer --- popolazione \(1{,}413\)B 2026 vs \(1{,}416\)B 2025
\item
  RAND Corporation 2025 --- fallimento politiche pronataliste cinesi
\item
  Crisis Group / AEI / INSS --- pressione PLA su Taiwan, \emph{Justice
  Mission 2025}
\item
  Lai 13/03/2025 --- \emph{foreign hostile force}
\item
  BRICS expansion 2025 (Egitto, Etiopia, Iran, EAU, Arabia Saudita,
  Indonesia) --- Wikipedia, RioTimesOnline
\item
  90\% trade RU-CN in valute locali, 50\% intra-BRICS in yuan ---
  InfoBRICS
\item
  Pilot token BRICS 31 ottobre 2025 --- BTCC, GIS Reports
\item
  BRICS Pay slittato a 2030; presidenza India 2026 con CBDC focus
\end{itemize}


